\subsection{连续性与等度连续性判别准则}

设 \(\mathbf{E}\) 与 \(\mathbf{F}\) 是同一除环 \(\mathbf{K}\) 上的拓扑向量空间；为使线性映射 \(f\)（\(\mathbf{E}\) 到 \(\mathbf{F}\) 中的映射）连续，只需它在原点连续 (GT, III, § 2.8, prop. 23)。该命题可推广如下：

命题 5. --- 设 \(E_{i}(1 \leqslant i \leqslant n)\) 与 F 是非离散赋值域 K 上的拓扑向量空间。为使 \(\prod_{i=1}^{n} E_{i}\) 到 F 中的多重线性映射 f 连

续于乘积空间 \(\prod_{i=1}^{n} E_i\) 中，只需它在 \((0, 0, ..., 0)\) 处连续。

设 \((a_{1}, a_{2}, \ldots, a_{n})\) 是 \(\prod_{i=1}^{n} \mathrm{E}_{i}\) 的任意一点；我们必须证明，对每个邻域 \(W\)（0 在 \(F\) 中的邻域），存在邻域 \(\mathrm{V}_{i}\)（0 在 \(\mathrm{E}_{i}\) 中的邻域）( \(1 \leqslant i \leqslant n\) )，使得诸关系式 \(z_{i} \in V\) 蕴含

\[
f \left(a _ {1} + z _ {1}, a _ {2} + z _ {2}, \dots , a _ {n} + z _ {n}\right) - f \left(a _ {1}, a _ {2}, \dots , a _ {n}\right) \in \mathrm{W}.
\]

现在，我们可以写成

\[
f (a _ {1} + z _ {1}, \dots , a _ {n} + z _ {n}) - f (a _ {1}, \dots , a _ {n}) = \sum_ {\mathrm{H}} u _ {\mathrm{H}}
\]

其中 H 取遍 \(2^{n}-1\) 个子集（即整数集 \(\{1,2,\ldots,n\}\) 的除 \(\{1,2,\ldots,n\}\) 本身以外的子集），而 \(u_{\mathrm{H}}=f(y_{1},y_{2},\ldots,y_{n})\)，其中 \(y_{i}=a_{i}\) 若 \(i\in H\)，而 \(y_{i}=z_{i}\) 若 \(i\notin H\)。存在 \(2^{n}-1\) 个 F 中 0 的均衡邻域 \(W_{H}\)，使得 \(\sum_{H}W_{H}\subset W\)；另一方面，由假设 f 在 \((0,0,\ldots,0)\) 处连续，故在每个 \(E_{i}\) 中存在 0 的一个邻域 \(U_{i}\) \(0(1\leqslant i\leqslant n)\)，使得 n 个关系式 \(x_{i}\in U_{i}\) 蕴含 \(f(x_{1},\ldots,x_{n})\in\bigcap_{\mathrm{H}}\mathrm{W}_{\mathrm{H}}\)。由于 \(U_{i}\) 是吸收的，存在 K 中的 \(\lambda_{i}\neq0\) 使得 \(\lambda_{i}a_{i}\in U_{i}\)。设 \(\lambda\) 是 K 中对每个子集 H 都满足 \(|\lambda|\geqslant\prod_{i\in H}|\lambda_{i}|^{-1}\) 的元素；我们证明邻域 \(V_{i}=\lambda^{-n}U_{i}\) 满足所需条件。可以写 \(u_{\mathrm{H}}=\mu f(x_{1},\ldots,x_{n})\)，其中 \(x_{i}\in U_{i}\)（对 \(1\leqslant i\leqslant n\)），而 \(\mu=\lambda^{-np}(\prod_{i\in H}\lambda_{i}^{-1})\)，p 是 \(\{1,2,\ldots,n\}\) 中不属于 H 的整数个数。由上式得 \(|\mu|\leqslant1\)，从而有 \(u_{H}\in\mu W_{H}\subset W_{H}\)，因为 \(W_{H}\) 均衡。命题得证。

命题 6. --- 在与命题 5 相同的关于 \(E_{i}(1 \leqslant i \leqslant n)\) 与 F 的假设下，为使 \(\prod_{i=1}^{n} E_{i}\) 到 F 中的多重线性映射所成的集 E 等度连续，只需该集在 \((0, 0, \ldots, 0)\) 处等度连续。

因为，在命题 5 的证明中，可以取 \(\mathrm{U}_i\) ( \(1 \leqslant i \leqslant n\) ) 使得关系式 \(x_i \in \mathrm{U}_i\) ( \(1 \leqslant i \leqslant n\) ) 蕴含 \(f(x_1, \ldots, x_n) \in \bigcap_{\mathrm{H}} \mathrm{W}_{\mathrm{H}}\) 对每个映射 \(f \in \mathcal{E}\) 成立。

